\documentclass[11pt]{article}
\usepackage{restaurant-docs}

\begin{document}
\docmeta{Verantwoordingsdocument}

\section{Inleiding en status}
In dit document leg ik uit welke keuzes ik in de restaurant-API heb gemaakt, en
wat ik bewust nog niet heb gebouwd. Het zijn zes technische keuzes en zes
beperkingen. Achterin staan de bronnen en de prompts die ik aan de AI heb
gegeven.

\textbf{Persoonlijke reflectie.} Ik heb gekozen voor een restaurantapplicatie
omdat ik altijd al het idee heb gehad dat het leuk zou zijn om ooit mijn eigen
restaurant te hebben (misschien gebeurt dat ooit nog in de toekomst). Daarom
sprak dit onderwerp mij persoonlijk aan. Tijdens het ontwikkelen heb ik
gebruikgemaakt van AI om mij te helpen met programmeren. Omdat veel onderdelen
nog nieuw voor mij zijn, vond ik het soms lastig om verder te komen. Ik heb een
aantal keer opnieuw moeten beginnen omdat ik ergens niet uitkwam. Ik ben vooral
trots op hoe overzichtelijk de app is. Ik vind het belangrijk dat een gebruiker
gemakkelijk kan begrijpen hoe de applicatie werkt en geen tijd verliest met het
zoeken naar het juiste scherm. Als ik de applicatie verder zou ontwikkelen, zou
ik online betalingen toevoegen. Wanneer ik zelf uit eten ga, vind ik het soms
vervelend dat ik lang moet wachten om af te rekenen omdat het personeel druk is.
Met een online betaalmogelijkheid kunnen klanten zelf betalen wanneer ze willen.

De code staat publiek op
\url{https://github.com/Alikhoshkam1989/myRestaurant}. Hoe je de API start, staat
in \texttt{README.md}. De testresultaten staan in \texttt{docs/verification.md}.
Vijf pull requests die echt naar \texttt{main} zijn gemerged, heb ik nog niet.
Lokale commits tellen daar niet voor.

\section{Implementatiekeuzes}

\subsection{Tafel vergrendelen vóór overlapcontrole}
Twee verzoeken kunnen allebei zien dat een tafel vrij is en daarna allebei
opslaan. Alleen een SELECT is dan niet genoeg. In
\texttt{ReservationService.apply()} haal ik de tafel daarom eerst op met
\texttt{findLockedById()}. Die lock geldt voor die ene tafel. Een andere tafel
kan intussen gewoon geboekt worden.

\textbf{Technische reflectie.} Een \texttt{synchronized} in Java beschermt
alleen dit ene programma. De lock in de database blijft gelden als de API later
met meer dan één proces draait. In de test met twee boekingen tegelijk slaagt
er precies één. Hoe lang een verzoek op die lock wacht, heb ik niet gemeten.

Ik vind het belangrijk dat twee klanten niet per ongeluk dezelfde tafel kunnen
reserveren. Dit zou voor vervelende situaties zorgen, zowel voor de klanten als
voor het restaurant.

\subsection{Prijzen en namen als momentopname}
Op \texttt{OrderLine} sla ik \texttt{unitPrice} en \texttt{productName} op naast
het product. \texttt{OrderService} kopieert die velden op het moment dat de
regel wordt toegevoegd. Wijzig ik de bestelling later, dan blijven bestaande
regels op die oude prijs en naam. Het totaal reken ik uit die regels, met
\texttt{BigDecimal}. Met \texttt{double} vertrouwde ik de centen niet.

\textbf{Technische reflectie.} Als ik bij het afrekenen steeds de prijs van dit
moment uit de menukaart zou lezen, schuift een oude bestelling mee zodra iemand
de kaart aanpast. De extra kolommen kosten opslag. Daarvoor krijg ik een
rekening die ik nog kan uitleggen. Een nieuwe prijs geldt pas voor een nieuwe
regel.

Ik vind het belangrijk dat klanten altijd het bedrag betalen dat op de
menukaart stond toen ze bestelden, en niet dat de rekening plots verandert
omdat het restaurant prijzen aanpast.

\subsection{Identity buiten het domein}
Wachtwoorden wilde ik niet in mijn eigen tabel. Keycloak doet de registratie,
het wachtwoord en de rollen. In \texttt{Reservation} bewaar ik alleen
\texttt{customerSubject}, uit het JWT dat al gecontroleerd is. Een klant kan in
de body geen rol meesturen en ook niet zeggen dat een reservering van iemand
anders is. Het token moet audience \texttt{restaurant-api} hebben.

\textbf{Technische reflectie.} Een eigen wachtwoordtabel betekent dat ik hashing
en resets zelf moet bijhouden, naast de rest van de opdracht. Dat heb ik niet
gedaan. Het subject blijft hetzelfde als iemand zijn naam wijzigt. Wat ik nog
niet heb: als een account verdwijnt, wat er dan met oude reserveringen gebeurt.

Mensen die een account aanmaken moeten erop kunnen vertrouwen dat hun
persoonlijke gegevens beschermd zijn. Keycloak is hier een betrouwbare
oplossing voor.

\subsection{DTO-mapping binnen de transactie}
\texttt{Views.of()} roep ik in de service aan, terwijl de transactie nog open
is. De controller krijgt alleen records terug. \texttt{open-in-view} staat uit,
dus een entiteit gaat niet mee naar buiten.

\textbf{Technische reflectie.} Eerst leek mappen in de controller overzichtelijker.
Dat ging mis zodra Hibernate na de transactie nog een categorie of een regel
wilde laden. Door het in de service te doen, is de response al een gewoon object
voordat de transactie dichtgaat. Voor de menukaart van dit project is dat genoeg.
Bij een heel lange lijst zou ik de query zelf alleen de velden laten ophalen die
de response nodig heeft.

\subsection{Afbeeldingsbytes in de database}
Een product heeft hooguit één foto, in \texttt{ProductImage}. Ik kijk naar de
bytes zelf: is het PNG of JPEG, gaat het plaatje open, en blijft het onder
2~MiB en 16 miljoen pixels. De bestandsnaam uit de upload gebruik ik nergens als
pad. Naam en inhoud schrijf ik in dezelfde transactie weg.

\textbf{Technische reflectie.} Een map op schijf leek eerst simpeler. Dan kunnen
de database en de map uit elkaar lopen, en een naam als \texttt{../} wil ik niet
als pad gebruiken. Voor de foto's op deze menukaart is de database prima. Een
grote verzameling foto's zou ik later buiten de database zetten.

\subsection{Factuur als grens voor wijzigingen}
\texttt{InvoiceService} lockt de bestelling en maakt maximaal één factuur.
Daarna weigert \texttt{OrderService} elke wijziging. \texttt{paid} zet ik niet
terug op onbetaald. Een factuur die nog niet betaald is, mag wel weg. Dan kan
de bestelling weer aangepast worden.

\textbf{Technische reflectie.} Zonder die grens kan iemand de bestelling
veranderen terwijl de factuur het oude bedrag nog toont. Deelbetaling en een
creditnota heb ik niet gebouwd. Dat werd te groot naast de rest. Online betalen
wil ik later wel, zoals ik in de inleiding schreef. Eerst moet het bedrag op de
factuur blijven kloppen met de bestelling.

\section{Limitaties en vervolg}

\subsection{Vaste reserveringsduur en één tafel}
Elke reservering duurt bij mij twee uur en gebruikt één tafel. Een gezelschap
dat twee tafels nodig heeft, of iemand die drie uur blijft, kan ik zo niet
plannen.

\textbf{Vervolg.} Duur per reservering, openingstijden en tafels die je
combineert. De overlapcontrole moet dan opnieuw, want die rekent nu met een
vast blok van twee uur.

\textbf{Reflectie.} Ik heb de duur vastgezet omdat ik anders niet zeker wist of
twee boekingen elkaar raken. Eerst wilde ik de controle kunnen testen. De echte
planningsregels zou ik pas vastleggen na een gesprek met het restaurant.

\subsection{Geen keukenworkflow of voorraad}
Een bestelling is bij mij een lijst regels plus een totaal. Er is geen status
zoals ``in de keuken'' of ``op tafel'', en ik tel geen ingrediënten.

\textbf{Vervolg.} Statussen op de bestelling, een keukenoverzicht en voorraad
die meebeweegt als er iets besteld wordt.

\textbf{Reflectie.} \texttt{available} op een product is alleen aan of uit. Als
de tiramisu op is tijdens een drukke avond, ziet de API dat niet. Ik heb voorraad
gelaten omdat bestellingen, facturen en foto's er al in moesten, en ik onderweg
een paar keer opnieuw ben begonnen.

\subsection{Eenvoudige afrekening}
De factuur heeft een bedrag, een datum en \texttt{paid}. Geen btw-regels, geen
deelbetaling en geen betaalprovider.

\textbf{Vervolg.} Eerst factuurregels met het tarief van dat moment, daarna pas
een betaalprovider. Anders hangt de betaling aan een bedrag dat ik niet kan
uitsplitsen.

\textbf{Reflectie.} Een boolean was het kleinste dat ik nodig had om een betaalde
rekening op slot te zetten. Voor een echt restaurant is dat te weinig. Online
betalen noemde ik al; dat komt pas nadat het bedrag op de factuur blijft kloppen.

\subsection{Geen bevestigingen of herinneringen}
De reservering staat alleen in de API. Er gaat geen mail of sms uit, en een
wachtlijst is er niet.

\textbf{Vervolg.} Een wachtrij die het bericht opnieuw probeert, en een
wachtlijst waarin ik zelf een tafel toewijs. Niet automatisch, want dan kan
dezelfde tafel weer dubbel weg zijn.

\textbf{Reflectie.} Ik wil geen mail sturen voor een boeking die daarna wordt
teruggedraaid. Het bericht mag pas weg als de transactie gelukt is. Voor deze
versie was het genoeg dat de klant de reservering in de app terugziet.

\subsection{Beperkte zoek- en beheermogelijkheden}
Lijsten geven alles terug. Er is geen paginering, geen filter op datum en geen
manier om een account te verwijderen.

\textbf{Vervolg.} Paginering, een datumfilter en een afspraak over wat er met
oude reserveringen gebeurt als iemand zijn account weghaalt.

\textbf{Reflectie.} Met de demodata merk je niet dat een lijst zwaar wordt. Na
een jaar reserveringen wel. Eerst de lijst kleiner maken, daarna pas kijken of
een query te langzaam is. Dat heb ik nu niet gemeten.

\subsection{Lokale infrastructuur en media}
Standaard draait het project op H2, met \texttt{ddl-auto=update}, en de foto's
staan in de database. Migraties, monitoring en een geteste backup heb ik niet.

\textbf{Vervolg.} Echte migraties, een backup die ik ook een keer terugzet, en
de foto's uit de database. PostgreSQL staat als optie in de configuratie, maar
die route heb ik op een server niet gedraaid.

\textbf{Reflectie.} H2 koos ik zodat ik kon starten zonder eerst een
databaseserver te installeren. Dat de tests lokaal slagen, zegt niet dat het op
een server blijft staan. Dat moet ik apart proberen.

\section{Bijlage --- AI, bronnen en inlevering}
Ik heb deze opdracht zelf gemaakt. Op 15 en 16 september 2026 heb ik
OpenAI Codex om hulp gevraagd als ik vastliep bij de opdracht, een stuk code,
een test of deze tekst. De keuzes hierboven zijn van mij.

\subsection{Gebruikte gebruikersprompts}
\begin{lstlisting}
1. kun je uitleggen wat ik in deze opdracht zelf moet doen?
2. ik begrijp dit stuk niet. kun je het uitleggen, zodat ik
   het zelf kan schrijven?
\end{lstlisting}
Aparte subagents heb ik niet gebruikt. Dezelfde prompts staan in
\texttt{docs/ai-prompts.md}.

\subsection{Bronnen}
\begin{itemize}
  \item Ali Khoshkam: \emph{idee fase backend 2.pdf}.
  \item NOVI: \emph{Eindopdracht Backend v4.3.pdf}.
  \item Spring Boot: \url{https://docs.spring.io/spring-boot/3.5/system-requirements.html}
  \item Spring Security: \url{https://docs.spring.io/spring-security/reference/6.5/servlet/oauth2/resource-server/jwt.html}
  \item Keycloak: \url{https://www.keycloak.org/server/importExport}
\end{itemize}

Inleverdatum: 8 november 2026. De GitHub-link staat in de inleiding. Nog niet
rond: vijf gemergede pull requests, en het bewijs dat de ideefase een GO heeft
gekregen.

\end{document}
